\documentclass[11pt]{article}
\usepackage[a4paper,margin=1in]{geometry}
\usepackage{xcolor}
\usepackage{titlesec}
\usepackage{enumitem}
\usepackage{hyperref}
\usepackage{parskip}
\usepackage{tcolorbox}

\definecolor{qcink}{HTML}{191919}
\definecolor{qcmuted}{HTML}{5a5955}
\definecolor{qccoral}{HTML}{d97757}
\definecolor{qcblue}{HTML}{6d8db5}
\definecolor{qcgreen}{HTML}{6b9c7c}
\definecolor{qcteal}{HTML}{2f7d70}
\definecolor{qcamber}{HTML}{f3c97a}
\definecolor{qcdark}{HTML}{2b2a24}
\definecolor{qccream}{HTML}{ece4d2}

\hypersetup{
    colorlinks=true,
    linkcolor=qcteal,
    urlcolor=qcteal,
    pdftitle={Quantum Channels --- Rules \& Concepts},
    pdfauthor={Waterloo Quantum Club},
}

\titleformat{\section}{\Large\bfseries\color{qcink}}{\thesection}{0.6em}{}
\titleformat{\subsection}{\large\bfseries\color{qcink}}{\thesubsection}{0.5em}{}

\newtcolorbox{rulebox}[1]{
    colback=qcdark,
    colframe=qcdark,
    coltext=qccream,
    arc=8pt,
    boxrule=0pt,
    left=12pt, right=12pt, top=8pt, bottom=8pt,
    title=\textcolor{qcamber}{\bfseries\small\MakeUppercase{#1}},
    fonttitle=\sffamily,
    coltitle=qcamber,
}

\newtcolorbox{conceptbox}[2]{
    colback=#1,
    coltext=qcink,
    colframe=#1,
    arc=8pt,
    boxrule=0pt,
    left=12pt, right=12pt, top=8pt, bottom=8pt,
    title=\textbf{\small\MakeUppercase{#2}},
    fonttitle=\sffamily,
    coltitle=qcink,
}

\begin{document}

\begin{center}
    {\Huge\bfseries Quantum Channels}\\[4pt]
    {\large\itshape A weekly strategy game inspired by quantum mechanics}\\[8pt]
    {\color{qcmuted} Waterloo Quantum Club}
\end{center}

\vspace{0.8em}
\hrule
\vspace{1.2em}

\section{What is this?}

\textbf{Quantum Channels} is a friendly weekly game you play in your browser. Every
week, you split 100 ``quanta'' across 10 channels and try to outscore your opponent.
The twist? Each week the rules change to mimic a real principle of quantum mechanics
--- so by playing, you build intuition for the strange behaviour of the quantum world.

You do \emph{not} need any physics background to play. The game teaches you as you go.

\section{How a round works}

\begin{itemize}[leftmargin=1.5em, itemsep=4pt]
    \item You have \textbf{100 quanta} to spread across \textbf{10 channels}.
    \item Slide each channel up or down to place quanta there. The total has to add up to 100 (or less).
    \item You can also place a \textbf{Superposition Token} on one channel (more on that below).
    \item Hit ``Submit'' to lock in your wavefunction.
\end{itemize}

\subsection*{When can I play?}

\begin{rulebox}{Weekly schedule}
\begin{itemize}[leftmargin=1.2em, itemsep=2pt]
    \item \textbf{Submit your wavefunction} --- any time, all week, until \textbf{Wednesday 9:00 AM}.
    \item \textbf{Resubmit any time} --- you can change your mind and update your wavefunction until the 9 AM deadline.
    \item \textbf{Submissions close} --- Wednesday 9 AM. Your opponent is revealed.
    \item \textbf{Results posted \& new week opens} --- Wednesday \textbf{7:00 PM}. The next week's rule is revealed, and submissions open again.
\end{itemize}
\end{rulebox}

\section{Scoring}

For each of the 10 channels, the player with \textbf{more quanta} in that channel wins it.
A tie gives no points. Each channel you win = \textbf{1 point}. Your weekly score is added
to your total on the leaderboard.

A perfect week is 10 points. A wipe-out is 0. Most games land in between.

\section{The superposition token}

Every week you get one Superposition Token. Tap any channel cell to place your token
there. The channel you choose matters: at resolution, the token only triggers if you
\textbf{win} that channel. So pick one you're confident about.

\begin{itemize}[leftmargin=1.5em]
    \item If you \textbf{lose} (or tie) the channel: the token does nothing.
    \item If you \textbf{win} the channel: a coin is flipped.
    \begin{itemize}
        \item \textbf{Heads (50\%)} --- the token grants \textbf{+3 bonus points}.
        \item \textbf{Tails (50\%)} --- the token grants \textbf{+0 points}.
    \end{itemize}
\end{itemize}

The token mirrors a real two-state quantum system: until you measure it, the
particle is in \emph{both} winning and losing states at once. Measurement
collapses it onto one outcome.

\section{Weekly rules}

This is the fun part. Each week introduces a new quantum rule, revealed only when the
new week opens. Each rule maps to a real concept from quantum mechanics. After the last
rule, the cycle repeats.

\begin{conceptbox}{qcblue!85}{Week 1 --- Quantum measurement}
\textbf{The rule:} when the week resolves, each player's \textbf{tallest channel
collapses to 0}.

\textbf{What it teaches:} in quantum mechanics, measuring a system disturbs it.
Just by looking at where a particle is, you change its state. The boldest bet in
your wavefunction is the one that gets ``measured'' first --- and that
measurement wipes it out.
\end{conceptbox}

\begin{conceptbox}{qccoral!85}{Week 2 --- Pauli exclusion}
\textbf{The rule:} for every channel where \emph{both} you and your opponent
placed quanta, \textbf{both stacks are wiped to 0}. The channel rejects everyone.

\textbf{What it teaches:} two identical fermions (like electrons) cannot occupy the
same quantum state at the same time --- the Pauli exclusion principle. In the game,
if you both try to ``occupy'' the same channel, neither of you gets it.
Predict where your opponent will bet --- and bet somewhere else.
\end{conceptbox}

\begin{conceptbox}{qcgreen!85}{Week 3 --- Quantum tunneling}
\textbf{The rule:} after both players submit, every channel rolls a 25\% chance to
\textbf{shift all of its quanta into the next channel to its right} (channel 10
wraps to channel 1). The channel left behind is empty.

\textbf{What it teaches:} quantum tunneling. A particle can sometimes leak through
a barrier that classical physics says is impossible. Your carefully-placed quanta
might just \emph{tunnel away}, leaving the original channel empty and stacking
into the next one.
\end{conceptbox}

\begin{conceptbox}{qccream}{Week 4 --- Heisenberg uncertainty}
\textbf{The rule:} after submission, one randomly-chosen channel per player has its
value \textbf{shifted by a random integer between $-15$ and $+15$ quanta} (the channel
can't go below 0). Scoring stays the same: 1 point per channel you win.

\textbf{What it teaches:} the uncertainty principle. Some pairs of properties (like
position and momentum) can't both be known exactly at the same time. The game adds a
dose of uncertainty to one of your channels --- you won't know which, or how much,
until resolution.
\end{conceptbox}

\section{Strategy tips}

\begin{itemize}[leftmargin=1.5em]
    \item \textbf{Spread, don't pile.} If the tallest-collapse rule is active, piling
          quanta on one channel guarantees you lose them. Even when it isn't, big
          stacks tend to be exposed by the weekly twist.
    \item \textbf{You only need +1.} You don't have to crush your opponent in a channel
          --- one extra quantum is enough to win it.
    \item \textbf{Use the token tactically.} The superposition token works best on a
          channel you don't mind losing. If it transfers, you've doubled up; if not,
          no harm done.
    \item \textbf{Read the rule of the week.} Don't play week 3 like week 1 --- the
          best distribution depends on what the universe is going to do to your
          quanta.
\end{itemize}

\section{Frequently asked}

\textbf{Can I see my opponent before the deadline?} \\
No. The opponent is hidden until Wednesday 9 AM, when submissions close.

\textbf{What if I forget to submit?} \\
You sit out the week --- no opponent, no score added. The physicist-bots only
appear when there's an \emph{odd} number of players who submitted (to keep the
matchmaking paired). So don't forget!

\textbf{Can I change my submission?} \\
Yes! You can resubmit as many times as you want until Wednesday 9 AM. Only the
last one counts.

\textbf{When does the new rule appear?} \\
The new week (and its rule) opens at Wednesday 7 PM, after the previous week's
results are posted.

\vspace{1.2em}
\hrule
\vspace{0.6em}
\begin{center}
    \small\color{qcmuted}
    Waterloo Quantum Club --- play at \href{https://uwquantum.com/game.html}{uwquantum.com/game.html}\\
    Questions? Find us on Discord.
\end{center}

\end{document}
